\documentclass[conference]{IEEEtran}

% -------------------------------------------------------
% Paquetes
% -------------------------------------------------------
\usepackage[spanish]{babel}
\usepackage[utf8]{inputenc}
\usepackage[T1]{fontenc}
\usepackage{amsmath}
\usepackage{amssymb}
\usepackage{graphicx}
\usepackage{booktabs}
\usepackage{hyperref}
\usepackage{cite}
\usepackage{pgfplots}
\pgfplotsset{compat=1.18}
\usepackage{tikz}
\usepackage{array}
\usepackage{multirow}
\usepackage{url}
\usepackage{fancyhdr}

% Nota: Asegúrese de tener el archivo 'uts_logonegro.png' en la carpeta
% de compilación (o ajuste la ruta relativa, ej. ../assets/uts_logonegro.png).

% -------------------------------------------------------
\begin{document}

\pagestyle{fancy}
\fancyhf{}
\rhead{\includegraphics[height=1.2cm]{uts_logonegro.png}}
\fancyfoot[C]{\thepage}
\renewcommand{\headrulewidth}{0.0pt}
\setlength{\headheight}{1.5cm}

% -------------------------------------------------------
% CABECERA
% -------------------------------------------------------
\title{Boletín | Centro de Inteligencia Artificial e Innovación\\
       \large Mes de [Mes] de [Año]}

\author{
  \IEEEauthorblockN{Autores: [Nombre del Autor 1]}
  \IEEEauthorblockN{[Nombre del Autor 2]}
  \IEEEauthorblockN{[Nombre del Autor 3]}
  \IEEEauthorblockA{Docentes: [Nombre Docente 1], [Nombre Docente 2]}
  \IEEEauthorblockA{
    Centro de Inteligencia Artificial e Innovación (CIAI)\\
    Colombia\\
    autor1@uts.edu.co\\
    autor2@uts.edu.co\\
    autor3@uts.edu.co
  }
}

\maketitle
\thispagestyle{fancy}

% -------------------------------------------------------
% RESUMEN
% -------------------------------------------------------
\begin{abstract}
Aquí va el texto del resumen general de esta edición del boletín del CIAI:
temas principales cubiertos durante el mes, lanzamientos destacados,
comparativas de rendimiento, herramientas y noticias relevantes.
\end{abstract}

\begin{IEEEkeywords}
palabra clave 1, palabra clave 2, palabra clave 3, CIAI, inteligencia artificial
\end{IEEEkeywords}

% ===============================================================
\section{Investigación}
% ===============================================================

Aquí va el texto introductorio de la sección de investigación: modelos,
publicaciones y avances analizados durante el periodo, junto con las
comparativas cuantitativas correspondientes.

% -------------------------------------------------------
\subsection{Nuevos Modelos}

% -- Organización / Proveedor 1 ------------------------------
\subsubsection{[Nombre de la Organización / Proveedor 1]}

Aquí va el texto de: descripción de los nuevos modelos lanzados, fecha de
publicación, capacidades técnicas, arquitectura y novedades principales,
con referencia de:~\cite{ejemplo_referencia}.

Aquí va el texto de: detalles adicionales sobre planes, actualizaciones de
ecosistema, seguridad o disponibilidad en API, con referencia
de:~\cite{ejemplo_referencia}.

% -- Organización / Proveedor 2 ------------------------------
\subsubsection{[Nombre de la Organización / Proveedor 2]}

Aquí va el texto de: descripción del segundo laboratorio o proveedor de IA,
avances reportados y contexto relevante, con referencia
de:~\cite{ejemplo_referencia}.

% -------------------------------------------------------
\subsection{Analítica de Rendimiento}

Aquí va el texto de: explicación de las métricas evaluadas (costo, latencia,
velocidad de generación, etc.) y fuente de los datos, con referencia
de:~\cite{ejemplo_referencia}.

\begin{table}[h]
\centering
\caption{Aquí va el título de la tabla comparativa -- [Mes Año]}
\label{tab:comparativa}
\begin{tabular}{@{}lcccc@{}}
\toprule
\textbf{Modelo} & \textbf{Métrica 1} & \textbf{Métrica 2} & \textbf{Métrica 3} & \textbf{Métrica 4} \\
                & (unidad)           & (unidad)           & (unidad)           & (unidad)           \\
\midrule
[Modelo 1] & 0.00 & 0.00 & 0.00 & 00 \\
[Modelo 2] & 0.00 & 0.00 & 0.00 & 00 \\
[Modelo 3] & 0.00 & 0.00 & 0.00 & 00 \\
[Modelo 4] & 0.00 & 0.00 & 0.00 & 00 \\
\bottomrule
\multicolumn{5}{@{}l}{\footnotesize *Aquí va una nota aclaratoria de la tabla (si aplica).} \\
\end{tabular}
\end{table}

\begin{figure}[h]
\centering
\begin{tikzpicture}
\begin{axis}[
  width=\columnwidth,
  height=5.5cm,
  ybar,
  symbolic x coords={
    Modelo 1,
    Modelo 2,
    Modelo 3,
    Modelo 4
  },
  xtick=data,
  ylabel={[Nombre de la Métrica 1 (unidad)]},
  ymin=0,
  bar width=0.35cm,
  nodes near coords,
  nodes near coords align={vertical},
  enlarge x limits=0.15,
  x tick label style={font=\tiny, rotate=30, anchor=east},
]
\addplot coordinates {
  (Modelo 1, 1.00)
  (Modelo 2, 0.50)
  (Modelo 3, 2.00)
  (Modelo 4, 1.50)
};
\end{axis}
\end{tikzpicture}
\caption{Aquí va la gráfica de: [descripción de la métrica 1 representada].
Fuente: [Nombre de la fuente], [Mes Año].}
\label{fig:grafica1}
\end{figure}

\begin{figure}[h]
\centering
\begin{tikzpicture}
\begin{axis}[
  width=\columnwidth,
  height=5.5cm,
  ybar,
  symbolic x coords={
    Modelo 1,
    Modelo 2,
    Modelo 3,
    Modelo 4
  },
  xtick=data,
  ylabel={[Nombre de la Métrica 2 (unidad)]},
  ymin=0,
  bar width=0.35cm,
  nodes near coords,
  nodes near coords align={vertical},
  enlarge x limits=0.15,
  x tick label style={font=\tiny, rotate=30, anchor=east},
]
\addplot coordinates {
  (Modelo 1, 1.25)
  (Modelo 2, 1.10)
  (Modelo 3, 3.40)
  (Modelo 4, 1.95)
};
\end{axis}
\end{tikzpicture}
\caption{Aquí va la gráfica de: [descripción de la métrica 2 representada].
Fuente: [Nombre de la fuente], [Mes Año].}
\label{fig:grafica2}
\end{figure}

% ===============================================================
\section{Herramientas}
% ===============================================================

Aquí va el texto introductorio de las herramientas, librerías o entornos de
desarrollo publicados o actualizados durante el mes.

% -------------------------------------------------------
\subsection{[Nombre de la Herramienta 1]}

Aquí va el texto de: descripción general de la herramienta, arquitectura,
funcionalidades principales y ventajas frente a otras alternativas, con
referencia de:~\cite{ejemplo_referencia}.

\begin{itemize}
  \item \textbf{Tipo:} [Aplicación web / desktop / librería]
  \item \textbf{Licencia:} [MIT / GPL-3.0 / Apache 2.0]
  \item \textbf{Web:} \url{https://ejemplo.com}
  \item \textbf{Repositorio:} \url{https://github.com/usuario/repositorio}
\end{itemize}

% -------------------------------------------------------
\subsection{[Nombre de la Herramienta 2]}

Aquí va el texto de: descripción de la segunda herramienta destacada y sus
características técnicas, con referencia de:~\cite{ejemplo_referencia}.

\begin{itemize}
  \item \textbf{Tipo:} [Tipo de herramienta]
  \item \textbf{Licencia:} [Tipo de licencia]
  \item \textbf{Repositorio:} \url{https://github.com/usuario/repositorio}
\end{itemize}

% ===============================================================
\section{Proyectos CIAI}
% ===============================================================

Aquí va el texto introductorio de los avances, publicaciones y actividades
del Centro de Inteligencia Artificial e Innovación durante el periodo.

% -------------------------------------------------------
\subsection{[Nombre del Proyecto 1]}

Aquí va el texto de: descripción del avance del proyecto, estado actual,
próximos pasos y autores o integrantes involucrados.

% ===============================================================
\section{Noticias}
% ===============================================================

Aquí va el texto introductorio de los eventos relevantes de la industria de la
Ingeniería de Sistemas, la inteligencia artificial y áreas afines.

% -------------------------------------------------------
\subsection{[Titular de la Noticia 1]}

Aquí va el texto de: desarrollo completo de la noticia, contexto, impacto y
datos clave, con referencia de:~\cite{ejemplo_referencia}.

\begin{itemize}
  \item \textbf{Fuente / Enlace oficial:} \url{https://ejemplo.com}
\end{itemize}

% -------------------------------------------------------
\subsection{[Titular de la Noticia 2]}

Aquí va el texto de: resumen y análisis de la segunda noticia relevante del
mes, con referencia de:~\cite{ejemplo_referencia}.

% -------------------------------------------------------
\bibliographystyle{IEEEtran}
\bibliography{referencias}

\end{document}
